\documentclass[11pt,a4paper]{article}
\usepackage[margin=1in]{geometry}
\usepackage{amsmath,amssymb,amsthm}
\usepackage{algorithm}
\usepackage{algpseudocode}
\usepackage{booktabs}
\usepackage{hyperref}

\theoremstyle{definition}
\newtheorem{definition}{\textbf{Definition}}[section]
\newtheorem{theorem}{\textbf{Theorem}}[section]
\newtheorem{lemma}[theorem]{\textbf{Lemma}}
\newtheorem{remark}{\textbf{Remark}}[section]

\title{\textbf{Model Merging: Task Arithmetic, TIES-Merging, DARE, and SLERP}}
\author{\textbf{Team Scaibu} \\ \small Email: \texttt{jaydeep.9664920749@gmail.com} \\ \small \textbf{Architecture and Core Engine Team}}
\date{\textbf{October 2026}}

\begin{document}
\maketitle

\section{Definitions and Cognitive Mental Models}

\subsection{Formal Mathematical Definition}
\textbf{Definition (Model Merging in Parameter Space):}
Model merging combines multiple neural network checkpoints fine-tuned from a common pretrained base model $W_{\text{base}} \in \mathbb{R}^d$ into a single consolidated checkpoint without requiring additional training data, GPU compute, or optimization steps.

For $M$ fine-tuned task models $\{W_m\}_{m=1}^M$, the \textbf{task vector} for model $m$ is defined as:
{\boldmath
\begin{equation}
\tau_m = W_m - W_{\text{base}}
\end{equation}
}
Different merging algorithms resolve interference and parameter conflicts across $\{\tau_m\}$:
\begin{enumerate}
    \item \textbf{Task Arithmetic} (Ilharco et al., 2023): Linear combination of task vectors:
    {\boldmath
    \begin{equation}
    W_{\text{merged}} = W_{\text{base}} + \sum_{m=1}^M \lambda_m \tau_m
    \end{equation}
    }
    \item \textbf{TIES-Merging (Trimming, Electing Sign, and Disjoint Merging)} (Yadav et al., 2023): Addresses parameter redundancy and interference via a three-stage pipeline:
    \begin{itemize}
        \item \emph{Trim}: Retain only the top-$k\%$ fraction of values in each $\tau_m$ by magnitude, setting the remaining values to zero: $\hat{\tau}_m = \operatorname{Trim}(\tau_m, k)$.
        \item \emph{Elect Sign}: Determine the consensus sign for each parameter index $j$:
        {\boldmath
        \begin{equation}
        s_j = \operatorname{sgn}\left(\sum_{m=1}^M \hat{\tau}_{m, j}\right)
        \end{equation}
        }
        \item \emph{Disjoint Merge}: For each parameter $j$, retain only the values that match sign $s_j$, and compute their mean:
        {\boldmath
        \begin{equation}
        \tau_{\text{TIES}, j} = \frac{1}{|A_j|} \sum_{m \in A_j} \hat{\tau}_{m, j}, \quad A_j = \{m \mid \operatorname{sgn}(\hat{\tau}_{m, j}) = s_j \text{ and } \hat{\tau}_{m, j} \neq 0\}
        \end{equation}
        }
    \end{itemize}
    \item \textbf{DARE (Drop And REscale)} (Yu et al., 2024): Sparsifies task vectors via Bernoulli dropout with probability $p \in [0, 1)$ and rescales surviving parameters by $\frac{1}{1-p}$ to preserve expectation:
    {\boldmath
    \begin{equation}
    \hat{\tau}_m = \frac{1}{1 - p} \left( M_{\text{mask}} \odot \tau_m \right), \quad M_{\text{mask}} \sim \operatorname{Bernoulli}(1 - p)
    \end{equation}
    }
    \item \textbf{SLERP (Spherical Linear Interpolation)}: Interpolates between two models along the hypersphere of parameter space, preserving weight vector norms:
    {\boldmath
    \begin{equation}
    W_{\text{SLERP}}(t) = \frac{\sin((1 - t)\theta)}{\sin \theta} W_1 + \frac{\sin(t\theta)}{\sin \theta} W_2, \quad \theta = \arccos\left(\frac{W_1^T W_2}{\|W_1\|_2 \|W_2\|_2}\right)
    \end{equation}
    }
\end{enumerate}

\subsection{Expanded Non-Technical Definition (Intuition for Anyone)}
\begin{enumerate}
    \item \textbf{Baking Multiple Recipes into One Cake}: You have one base model. Chef A fine-tuned it to be a master coder; Chef B fine-tuned it to be a brilliant doctor. Model merging blends both brains together in 5 seconds without burning any GPU electricity!
    \item \textbf{Eliminating Cognitive Interference (TIES)}: When two experts disagree on a specific parameter, simply averaging their numbers causes destructive interference (canceling out). TIES identifies which expert is confident, discards quiet noise, elects a majority sign, and averages only agreeing values.
    \item \textbf{Extreme Pruning (DARE)}: Up to 90\% of fine-tuned weight adjustments are completely redundant! DARE randomly zeroes out 90\% of the delta and boosts the remaining 10\%, making multiple task vectors orthogonal and conflict-free.
    \item \textbf{Geometric Blending (SLERP)}: Instead of cutting through the dense center of high-dimensional space where neural weights lose their natural magnitude, SLERP glides smoothly along the curved sphere of the weight manifold.
\end{enumerate}

\subsection{Child-Friendly Mental Model (Physical Metaphor)}
Imagine three sculptors working on copies of the same clay statue:
\begin{itemize}
    \item \textbf{Sculptor 1} carves wings ($\tau_1$).
    \item \textbf{Sculptor 2} carves armor ($\tau_2$).
    \item \textbf{Simple Addition}: Gluing both changes together might cause wings and armor to collide and smash into dust.
    \item \textbf{TIES-Merging}: Examining the clay changes, shaving away minor thumbprints (trimming), resolving collisions by voting (electing sign), and smoothly combining the wings and armor into one glorious statue!
\end{itemize}

\section{Core Mathematical Equations and Definitions}

\subsection{Task Vector Formulation}
{\boldmath
\begin{equation}
\tau = W_{\text{task}} - W_{\text{base}}
\end{equation}
}
\begin{itemize}
    \item \textbf{Formal Definition}: The displacement vector in parameter space spanning the base model to the fine-tuned checkpoint.
    \item \textbf{What It Means}: Isolates the learned task expertise from the foundation representations.
    \item \textbf{Why It Is Important}: Allows algebraic manipulation of capabilities (addition, scaling, subtraction/unlearning).
\end{itemize}

\subsection{TIES Sign Election}
{\boldmath
\begin{equation}
s_j = \operatorname{sgn}\left(\sum_{m=1}^M \hat{\tau}_{m, j}\right)
\end{equation}
}
\begin{itemize}
    \item \textbf{Formal Definition}: Majority direction vote over trimmed parameter values at coordinate $j$.
    \item \textbf{What It Means}: Discards minority conflicting parameter shifts that cause destructive cancellation.
    \item \textbf{Why It Is Important}: Eliminates cross-model parameter cancellation when merging many specialized models.
\end{itemize}

\subsection{SLERP Normalized Geodesic Step}
{\boldmath
\begin{equation}
W(t) = \frac{\sin((1 - t)\theta)}{\sin \theta} W_1 + \frac{\sin(t\theta)}{\sin \theta} W_2
\end{equation}
}
\begin{itemize}
    \item \textbf{Formal Definition}: Constant-speed shortest-path interpolation along the surface of a high-dimensional unit hypersphere.
    \item \textbf{What It Means}: Rotates smoothly between two model orientations without reducing weight norm length.
    \item \textbf{Why It Is Important}: Avoids activation scale collapse caused by standard Euclidean linear averaging.
\end{itemize}

\section{Rigorous Mathematical Analysis Checklist}

\begin{itemize}
    \item \textbf{Notation}: $W_{\text{base}} \in \mathbb{R}^d$, $\{W_m\}_{m=1}^M$ (task checkpoints), $\tau_m = W_m - W_{\text{base}}$ (task vectors), $\lambda_m$ (merge coefficients).
    \item \textbf{Epistemic Status}: Ilharco et al. (ICLR 2023), Yadav et al. (NeurIPS 2023), Yu et al. (ICML 2024). Standard state-of-the-art methodology for zero-cost multi-task synthesis.
    \item \textbf{Dependency Graph}: Checkpoints $\{W_m\} \to \text{Task Deltas } \{\tau_m\} \to \text{Trim / Drop / Sign Election} \to \text{Aggregate } \tau_{\text{merged}} \to W_{\text{base}} + \tau_{\text{merged}}$.
    \item \textbf{Dimensions}: Total flattened parameter count $d$ across all model layers.
    \item \textbf{Regimes}: Multi-domain model unification, safety alignment transfer, task unlearning/censorship removal.
    \item \textbf{Invariants}: If $\tau_m = \mathbf{0}$, the merged model is identically equal to $W_{\text{base}}$.
    \item \textbf{Isomorphisms}: Isomorphic to manifold Riemannian geodesic averaging and sparse subspace voting.
\end{itemize}

\section{Theoretical Theorems and Formal Proofs}

\begin{theorem}[Unbiased Expectation of DARE Sparsified Task Vectors]
Let $\tau \in \mathbb{R}^d$ be an arbitrary fixed task vector. For dropout probability $p \in [0, 1)$, let mask entries $m_j \sim \operatorname{Bernoulli}(1 - p)$ be independent identically distributed binary random variables, and define the DARE vector:
{\boldmath
\begin{equation}
\hat{\tau}_j = \frac{m_j}{1 - p} \tau_j
\end{equation}
}
Then the expected value of $\hat{\tau}$ is strictly equal to $\tau$:
{\boldmath
\begin{equation}
\mathbb{E}[\hat{\tau}] = \tau
\end{equation}
}
and the variance of entry $j$ is $\operatorname{Var}(\hat{\tau}_j) = \frac{p}{1 - p} \tau_j^2$.
\end{theorem}

\begin{proof}
By definition of the Bernoulli distribution with parameter $q = 1 - p$:
{\boldmath
\begin{equation}
\mathbb{E}[m_j] = 1 - p, \quad \operatorname{Var}(m_j) = (1 - p)(1 - (1 - p)) = p(1 - p)
\end{equation}
}
Taking expectation of the $j$-th entry of $\hat{\tau}$:
{\boldmath
\begin{equation}
\mathbb{E}[\hat{\tau}_j] = \mathbb{E}\left[\frac{m_j}{1 - p} \tau_j\right] = \frac{\tau_j}{1 - p} \mathbb{E}[m_j] = \frac{\tau_j}{1 - p} (1 - p) = \tau_j
\end{equation}
}
Thus, $\mathbb{E}[\hat{\tau}] = \tau$, establishing that DARE is an unbiased estimator of the original task vector.

Now computing the variance:
{\boldmath
\begin{equation}
\operatorname{Var}(\hat{\tau}_j) = \operatorname{Var}\left(\frac{m_j}{1 - p} \tau_j\right) = \frac{\tau_j^2}{(1 - p)^2} \operatorname{Var}(m_j) = \frac{\tau_j^2}{(1 - p)^2} p(1 - p) = \frac{p}{1 - p} \tau_j^2
\end{equation}
}
This proves that DARE preserves the exact expected directional update while introducing extreme parameter sparsity ($p \approx 0.9$).
\end{proof}

\begin{theorem}[Preservation of Hypersphere Geodesic Length in SLERP]
Let $u, v \in \mathbb{R}^d$ be two unit vectors ($\|u\|_2 = \|v\|_2 = 1$) with angle $\theta = \arccos(u^T v) \in (0, \pi)$. For any interpolation parameter $t \in [0, 1]$, the SLERP vector:
{\boldmath
\begin{equation}
w(t) = \frac{\sin((1 - t)\theta)}{\sin \theta} u + \frac{\sin(t\theta)}{\sin \theta} v
\end{equation}
}
is a unit vector ($\|w(t)\|_2 = 1$) and lies on the great circle connecting $u$ and $v$ with angle to $u$ equal to $t\theta$.
\end{theorem}

\begin{proof}
Compute the squared Euclidean norm $\|w(t)\|_2^2 = w(t)^T w(t)$:
{\boldmath
\begin{align}
\|w(t)\|_2^2 &= \frac{1}{\sin^2 \theta} \left[ \sin^2((1 - t)\theta) u^T u + \sin^2(t\theta) v^T v + 2 \sin((1 - t)\theta) \sin(t\theta) u^T v \right]
\end{align}
}
Since $\|u\|_2 = \|v\|_2 = 1$ and $u^T v = \cos \theta$:
{\boldmath
\begin{align}
\|w(t)\|_2^2 &= \frac{1}{\sin^2 \theta} \left[ \sin^2((1 - t)\theta) + \sin^2(t\theta) + 2 \sin((1 - t)\theta) \sin(t\theta) \cos \theta \right]
\end{align}
}
Using the trigonometric identity $\cos(A - B) = \cos A \cos B + \sin A \sin B$ and angle sum formulas:
Let $A = (1 - t)\theta$ and $B = t\theta$, then $A + B = \theta$.
Notice that:
{\boldmath
\begin{align}
\sin^2 A + \sin^2 B + 2 \sin A \sin B \cos(A + B) &= \sin^2 A + \sin^2 B + 2 \sin A \sin B (\cos A \cos B - \sin A \sin B) \\
&= \sin^2 A + \sin^2 B + 2 \sin A \cos A \sin B \cos B - 2 \sin^2 A \sin^2 B \\
&= \sin^2 A (1 - \sin^2 B) + \sin^2 B (1 - \sin^2 A) + 2 \sin A \cos A \sin B \cos B \\
&= \sin^2 A \cos^2 B + \cos^2 A \sin^2 B + 2 (\sin A \cos B)(\cos A \sin B) \\
&= (\sin A \cos B + \cos A \sin B)^2 = \sin^2(A + B) = \sin^2 \theta
\end{align}
}
Substituting this numerator back into the norm:
{\boldmath
\begin{equation}
\|w(t)\|_2^2 = \frac{\sin^2 \theta}{\sin^2 \theta} = 1 \implies \|w(t)\|_2 = 1
\end{equation}
}
Thus, $w(t)$ strictly traverses the unit hypersphere surface without radial distortion.
\end{proof}

\section{Foundational Architectural Questions and Epistemic Meta-Questions}

\subsection{Architectural Question 1: Unlearning via Task Vector Negation}
\textbf{Question:} Can model merging be used to remove unwanted capabilities (e.g. copyright, toxicity, cyber-exploit knowledge)?\\
\textbf{Answer:} Yes. In Task Arithmetic, subtracting a task vector fine-tuned on undesirable data removes that capability:
{\boldmath
\begin{equation}
W_{\text{safe}} = W_{\text{base}} - \lambda \tau_{\text{toxic}}
\end{equation}
}
Empirical studies confirm that setting $\lambda \in [0.5, 1.0]$ effectively degrades toxic generation while retaining over 98\% of the model's general knowledge and conversational fluency.

\textbf{Meta-Question:} Why does simple linear averaging fail when merging many models?\\
\textbf{Answer:} When merging 5 or more models, parameters at identical coordinates frequently have opposite signs (some positive, some negative). Simple averaging sums these to near zero, destroying the fine-tuned capabilities of all models. TIES resolves this by electing the majority sign.

\subsection{Architectural Question 2: DARE with TIES Pipeline (DARE-TIES)}
\textbf{Question:} What is the state-of-the-art merging workflow used in open-source leaderboards (e.g. Open LLM Leaderboard)?\\
\textbf{Answer:} The highest-performing merge recipe is DARE-TIES:
1. Apply DARE to drop $80\%\text{--}90\%$ of parameters in each task vector and rescale the rest.
2. Apply TIES sign election across the sparsified task vectors to eliminate remaining sign conflicts.
3. Compute the disjoint average and add back to $W_{\text{base}}$.

\textbf{Meta-Question:} Does model merging require fine-tuned models to share the exact same base model?\\
\textbf{Answer:} Yes. The checkpoints must share an identical base weight initialization and token vocabulary. Merging checkpoints derived from different base architectures (e.g. LLaMA-3 and Mistral) requires cross-model representation alignment or weight permutation matching (Git Re-Basin).

\section{Algorithmic Implementation Specification}

\begin{algorithm}[H]
\caption{Model Merging via TIES and SLERP}
\begin{algorithmic}[1]
\Require Base weights $\mathbf{W}_{\text{base}} \in \mathbb{R}^d$, task models $\{\mathbf{W}_m\}_{m=1}^M \in \mathbb{R}^{M \times d}$, merge method $\in \{\text{task\_arithmetic}, \text{ties}, \text{slerp}\}$, top fraction $k_f$, coefficients $\mathbf{\lambda} \in \mathbb{R}^M$
\Ensure Merged weights $\mathbf{W}_{\text{merged}} \in \mathbb{R}^d$, consensus count $C$
\State Compute task vectors $\mathbf{\tau}_m \gets \mathbf{W}_m - \mathbf{W}_{\text{base}}$ for $m \in \{1, \dots, M\}$
\If{$\text{method} == \text{"ties"}$}
    \For{$m = 1$ \textbf{to} $M$}
        \State Retain top $\lceil k_f \cdot d \rceil$ entries of $\mathbf{\tau}_m$ by magnitude; set rest to $0$
    \EndFor
    \For{$j = 0$ \textbf{to} $d-1$}
        \State $s_j \gets \operatorname{sgn}\left(\sum_{m=1}^M \mathbf{\tau}_{m}[j]\right)$
        \State Gather agreeing entries $A_j \gets \{ \mathbf{\tau}_m[j] \mid \operatorname{sgn}(\mathbf{\tau}_m[j]) == s_j \text{ and } \mathbf{\tau}_m[j] \neq 0 \}$
        \If{$|A_j| > 0$}
            \State $\mathbf{W}_{\text{merged}}[j] \gets \mathbf{W}_{\text{base}}[j] + \frac{1}{|A_j|} \sum_{v \in A_j} v$
        \Else
            \State $\mathbf{W}_{\text{merged}}[j] \gets \mathbf{W}_{\text{base}}[j]$
        \EndIf
    \EndFor
\ElsIf{$\text{method} == \text{"slerp"}$}
    \State Compute angle $\theta \gets \arccos\left(\frac{\mathbf{W}_1^T \mathbf{W}_2}{\|\mathbf{W}_1\|_2 \|\mathbf{W}_2\|_2}\right)$
    \State $\mathbf{W}_{\text{merged}} \gets \frac{\sin((1 - t)\theta)}{\sin \theta} \mathbf{W}_1 + \frac{\sin(t\theta)}{\sin \theta} \mathbf{W}_2$
\Else
    \State $\mathbf{W}_{\text{merged}} \gets \mathbf{W}_{\text{base}} + \sum_{m=1}^M \mathbf{\lambda}_m \mathbf{\tau}_m$
\EndIf
\State \Return $(\mathbf{W}_{\text{merged}}, C)$
\end{algorithmic}
\end{algorithm}

\section{Big-$\Theta$ Complexity Analysis}

\begin{itemize}
    \item \textbf{Task Vector Extraction Time}: $\Theta(M \cdot d)$ subtractions.
    \item \textbf{TIES Trimming Time}: $\Theta(M \cdot d \log d)$ for sorting parameter magnitudes (or $\Theta(M \cdot d)$ with quickselect).
    \item \textbf{Sign Election and Merging Time}: $\Theta(M \cdot d)$ operations.
    \item \textbf{SLERP Interpolation Time}: $\Theta(d)$ operations for dot product, norms, and linear combination.
    \item \textbf{Memory Storage Overhead}: $\Theta(M \cdot d)$ floats in system RAM during merge calculation.
\end{itemize}

\section{Single Canonical Repository Reference}

The complete deterministic scalar implementation, verification suite, and unit test benchmarks are published at:
\begin{center}
\url{https://github.com/Chief-Strategist-J/llm-obs-infra}
\end{center}

\end{document}
